% ============================================================
% Capitulo 2 -- Fundamentacao Teorica
% Tamanho-alvo: 12-18 paginas
% Ordem de escrita: PENULTIMO (depois de Resultados -- voce sabe melhor
% o que precisa fundamentar quando ja viu os dados).
% Sem figuras neste capitulo (toda visualizacao vai pra Resultados).
% ============================================================

\chapter{Fundamenta\c{c}\~ao Te\'orica}
\label{cap:fundamentacao}

% ------------------------------------------------------------
\section{Detec\c{c}\~ao de fake news: panorama}
\label{sec:panorama}

% ESCREVER (1.5 pagina): NLP classico, contexto social, propagacao como
% sinal complementar.
% Citacoes: \cite{pmc_overviewfakenews}, \cite{ieee_fakenewslandscape},
% \cite{mdpi_cnn_llm_nlp}.
% Tese desta secao: "abordagens isoladas (NLP-only ou estrutura-only) tem
% limites; trabalhos recentes movem para hibridizacao".

% ------------------------------------------------------------
\section{Grafos e redes sociais}
\label{sec:grafos}

% ESCREVER (1 pagina): definicoes basicas (grafo dirigido, arvore de
% propagacao, raiz, profundidade, branching factor). Sem citacao
% obrigatoria; pode usar \cite{newman2003structure} se quiser textbook.
% Topico-ancora: propagacao numa rede social como cascata enraizada
% (post original = raiz; reposts/replies = filhos diretos ou indiretos).

% ------------------------------------------------------------
\section{Embeddings textuais}
\label{sec:embeddings}

% ESCREVER (1 pagina): word embeddings -> BERT -> Sentence-BERT.
% Citacoes obrigatorias: \cite{devlin2019bert}, \cite{reimers2019sbert}.
% Modelo especifico: paraphrase-multilingual-mpnet-base-v2 (768 dim).
% Pegadinha: Sentence-BERT NAO e BERT vanilla. Nao diga "usamos BERT" --
% diga "usamos Sentence-BERT".

% ------------------------------------------------------------
\section{Redes Neurais de Grafos}
\label{sec:gnns}

% ESCREVER (3-4 paginas): subsecoes 2.4.1, 2.4.2, 2.4.3 (uma por arquitetura).

\subsection{Graph Convolutional Networks (GCN)}
\label{sec:gcn}

% \cite{kipf2017gcn}. Equacao central, intuicao, complexidade.
% Regra de propagacao:
\begin{equation}
H^{(l+1)} = \sigma\left(\hat{D}^{-1/2}\hat{A}\hat{D}^{-1/2} H^{(l)} W^{(l)}\right)
\label{eq:gcn}
\end{equation}
% onde $\hat{A} = A + I$ (auto-loops) e $\hat{D}$ e a matriz diagonal de
% graus de $\hat{A}$.

\subsection{Graph Attention Networks (GAT)}
\label{sec:gat}

% \cite{velickovic2018gat}. Introduz coeficiente de atencao por aresta:
\begin{equation}
\alpha_{ij} = \frac{\exp\left(\mathrm{LeakyReLU}(\mathbf{a}^T [W h_i \,\|\, W h_j])\right)}{\sum_{k \in \mathcal{N}(i)} \exp\left(\mathrm{LeakyReLU}(\mathbf{a}^T [W h_i \,\|\, W h_k])\right)}
\label{eq:gat_alpha}
\end{equation}

\subsection{GraphSAGE}
\label{sec:sage}

% \cite{hamilton2017graphsage}. Agregacao indutiva: amostra vizinhanca + agrega.
% Mencionar: "Implementamos com PyTorch Geometric \cite{fey2019pyg}".
% Mencionar: pooling final = global_mean_pool (alternativas: add, max).

% ------------------------------------------------------------
\section{Degenera\c{c}\~ao do GCN com features homog\^eneas}
\label{sec:degeneracao}

% SECAO CRITICA -- promessa explicita do outline v2.2. Conteudo obrigatorio:
%
% 1. Notacao: $X \in \mathbb{R}^{N \times d}$, $\tilde{A} = \hat{D}^{-1/2}\hat{A}\hat{D}^{-1/2}$.
% 2. Cenario: suponha $X_i = X_j \, \forall i, j$. Entao $X = \mathbf{1}_N x^T$.
% 3. Demonstracao matematica (abaixo).
% 4. Conclusao: $h_i^{(1)}$ e multiplo escalar do mesmo $xW$, com fator
%    dependente apenas do grau de $i$.
% 5. Implicacao: justifica adicionar features posicionais (\ref{sec:encodings}).

% Demonstracao a desenvolver na prosa:
\begin{align}
H^{(1)} &= \sigma\left(\tilde{A} X W\right) \notag \\
        &= \sigma\left(\tilde{A} \mathbf{1}_N x^T W\right) \notag \\
        &= \sigma\left((\tilde{A} \mathbf{1}_N)(x W)^T\right)
\label{eq:degeneracao}
\end{align}

% O termo $(\tilde{A} \mathbf{1}_N)$ e vetor de "grau normalizado" -- mesmo
% escalar para cada no, modulado apenas pelo seu grau.
% Sem informacao para discriminar entre nos da mesma vizinhanca.

% ------------------------------------------------------------
\section{Encodings posicionais}
\label{sec:encodings}

% ESCREVER (1 pagina): definir as 3 features adicionadas.
% - is_root \in \{0, 1\}: 1 se no e raiz da arvore.
% - grau_norm = deg(v) / deg_max_global \in [0, 1].
% - pos = ordem BFS do no / num_filhos, normalizado em [0, 1].
% Justificar: quebram a degenerencia de \S\ref{sec:degeneracao}.
% Citacao opcional relacionada: \cite{arxiv_topologicalfeatures}.

% ------------------------------------------------------------
\section{Explicabilidade em GNNs: GNNExplainer}
\label{sec:gnn_explainer}

% ESCREVER (1 pagina): mascara de arestas, mascara de features, problema
% de otimizacao (mutual information).
% Citacao obrigatoria: \cite{ying2019gnnexplainer}.
% Mencionar uso pratico: torch_geometric.explain.Explainer com
% algorithm=GNNExplainer(epochs=200), edge_mask_type="object",
% modo multiclass_classification em nivel de grafo.

% ------------------------------------------------------------
\section{Valida\c{c}\~ao estat\'istica}
\label{sec:estatistica}

% ESCREVER (1 pagina): k-fold estratificado, teste t pareado vs independente
% (justificar o pareado), Cohen's d.
% Citacao obrigatoria: \cite{cohen1988statistical}.
%
% Tabela inline de interpretacao de Cohen's $d$ (Cohen 1988):

\begin{table}[h]\centering
\caption{Interpreta\c{c}\~ao de tamanho de efeito de Cohen ($d$).}
\label{tab:cohen_d_interpretacao}
\begin{tabular}{ll}
\toprule
$|d|$           & Interpreta\c{c}\~ao \\
\midrule
$< 0.2$         & desprez\'ivel \\
$0.2 \leq |d| < 0.5$ & pequeno \\
$0.5 \leq |d| < 0.8$ & m\'edio \\
$\geq 0.8$      & grande \\
\bottomrule
\end{tabular}
\end{table}

% ------------------------------------------------------------
\section{Trabalhos relacionados}
\label{sec:relacionados}

% ESCREVER (2-3 paginas):

\subsection{Surveys de GNNs em fake news}
\label{sec:surveys}
% \cite{pmc_gnnsurveyfakenews}, \cite{pmc_overviewfakenews}.

\subsection{M\'etodos GNN can\^onicos}
\label{sec:metodos_canonicos}
% \cite{monti2019fakenews}, \cite{dou2021upfd}, \cite{bian2020bigcn},
% \cite{mdpi_bidirectionalgcn}.

\subsection{Detec\c{c}\~ao multil\'ingue}
\label{sec:multilingue}
% \cite{acl_hemt}, \cite{crosslanguagefakenews}, \cite{pmc_multilangs_covid}.

\subsection{Trabalhos brasileiros relacionados}
\label{sec:trabalhos_br}
% \cite{ufc_fakewhatsapp}, \cite{pucsp_gnnfakenews}, \cite{usp_fakenews}.

\subsection{Homofilia e assortatividade em redes}
\label{sec:homofilia}
% Fundamentacao teorica para por que topologia funciona em alguns dominios
% e nao em outros. Citacoes: \cite{newman2003mixing}, \cite{pei2020geomgcn}.

\subsection{Posicionamento deste trabalho}
\label{sec:posicionamento}

% Texto sugerido (paragrafo final):
% "O presente trabalho NAO propoe uma nova arquitetura GNN. Sua contribuicao
% e uma analise sistematica + diagnostico sobre quando o paradigma topologico
% funciona, ancorando achados positivos e negativos em uma tese estrutural
% unificada. Limitacao assumida: nao rodamos BiGCN \cite{bian2020bigcn} ou
% GCNFN diretamente -- comparacao numerica com esses metodos e feita apenas
% via valores publicados (ver \S\ref{sec:limitacoes})."
